\documentclass[landscape, 8pt]{extarticle}

\usepackage{../../preamble}
\usepackage[separate]{../../rss/thmboxes_v5}

%For \upgamma symbol for graphs
\usepackage{upgreek}


\begin{document}
	
	\fontsize{5pt}{6pt} \selectfont
	
	%Reduce horizontal padding in 'tabular' environment, make rows taller for legibility
	\setlength\tabcolsep{4.5pt}
	\setlength\extrarowheight{2pt}
	
	\setlength{\abovedisplayskip}{0pt}
	\setlength{\belowdisplayskip}{0pt}
	\setlength{\abovedisplayshortskip}{0pt}
	\setlength{\belowdisplayshortskip}{0pt}
	
	\begin{multicols*}{4}
		\raggedcolumns 
		%Original template by Leon, modified by Louis. Content by Louis
		\vspace{-5pt}
		\subsection*{HAna Exam Sheet 2025/26 - Back From The Dead}
		\vspace{-5pt}
		
		
		
		\begin{rem}[Notation]{rem:notati}{ }
			\begin{itemize}
				\item $\mathbb{N}_0 = \mathbb{N} \cup \{0\}$
			\end{itemize}
			
		\end{rem}
		
		
		\begin{thm}[Random Identities]{thm:randints}{ }
			\begin{itemize}
				\item $\sum_{n=1}^{\infty} 2^{-n}\epsilon =\epsilon$ for $\epsilon > 0$.
			\end{itemize}
			
		\end{thm}
		
		
	%--------------------------------------------------------------------------------------%
	%%%%%%%%%%%%%%%%%%%%%%%%%%%%%%%%%%%%% SECTION 1 %%%%%%%%%%%%%%%%%%%%%%%%%%%%%%%%%%%%%%%%
	%--------------------------------------------------------------------------------------%
	
	
	{\footnotesize \textbf{\underline{Reals, Sequences, and Series}}}
	
	
	{\tiny \textbf{\underline{Real Sequences}}}
		
		
		
		\begin{dfn}[Bounds and Boundedness]{dfn:bounds}{ }
			Let $A \subseteq R$.
			\begin{itemize}[label=*]
				\item $x \in \mathbb{R}$ is an \textbf{upper bound for $A$} if $x \ge a$ for all $a \in A$
				\item $A$ is \textbf{bounded above} if there exists an upper bound $x \in \mathbb{R}$ for $A$.
				\item $s \in \mathbb{R}$ is a \textbf{least upper bound} or \textbf{supremum} for $A$ if $s$ is an upper bound for $A$, and $x \ge s$ whenever $x$ is an upper bound for $A$.
				\item \textbf{Lower bounds} stuff is akin to this - if u don't know it go back to FPM
				\item $A$ is \textbf{bounded} if it is bounded above and below
			\end{itemize}
			
		\end{dfn}
		
		\begin{axi}[The Completeness Axiom]{axi:test}{ }
			Every bounded above non-empty $A \subseteq \mathbb{R}$ has a least upper bound in $\mathbb{R}$
			
		\end{axi}
		
		
		\begin{dfn}[Nested Sets]{dfn:nestsets}{ }
			A sequence $(I_n)_{n \in \mathbb{N}}$ of sets is \textbf{nested} if $I_n \supseteq I_{n+1}$ for all $n \in \mathbb{N}$
			
		\end{dfn}
		
		
		\begin{thm}[The Nested Interval Theorem]{thm:nestints}{ }
			If $(I_n)_{n \in \mathbb{N}}$ is a nested sequence of nonempty closed bounded intervals, then the set $E := \bigcap_{n \in \mathbb{N}}I_n = \left\{x \in \mathbb{R}: x \in I_n \forall n \in \mathbb{N}\right\}$ is nonempty. \\
			If $\lambda(I_n) \to 0$ as $n \to \infty$ then $E$ contains exactly one element.
			
		\end{thm}
		
		
		\begin{dfn}[Convergence]{dfn:cnvgnce}{ }
			A real sequence $(a_n)_{n \in \mathbb{N}}$ \textbf{converges} to $a \in \mathbb{R}$ iff for every $\epsilon > 0$ there exists $N \in \mathbb{N}$ such that $\left|a_n - a\right| <\epsilon$ for all $n \ge N$. A sequence which does not converge to some real number \textbf{diverges}.
			\begin{itemize}
				\item We say $a_n \to \infty$ as $n \to \infty$ if $\forall R \in \mathbb{R}$, $\exists N \in \mathbb{N}$ s.t. $a_n > R$ for all $n \ge \mathbb{N}$
				\item We say $a_n \to -\infty$ as $n \to \infty$ if $\forall R \in \mathbb{R}$, $\exists N \in \mathbb{N}$ s.t. $a_n < R$ for all $n \ge \mathbb{N}$
			\end{itemize}
			
		\end{dfn}
		
		
		\begin{dfn}[Monotonicity]{dfn:monton}{ }
			Let $(a_n)_{n \in \mathbb{N}}$ be a real sequence and $A := \{a_n : n \in\mathbb{N}\}$. Then:
			\begin{enumerate}[label=(\arabic*)]
				\item $(a_n)_{n \in \mathbb{N}}$ is \textbf{monotone non-decreasing} if $a_{n + 1} \ge a_n$ for all $n \in \mathbb{N}$
				\item $(a_n)_{n \in \mathbb{N}}$ is \textbf{monotone non-increasing} if $a_{n + 1} \le a_n$ for all $n \in \mathbb{N}$
				\item $(a_n)_{n\in \mathbb{N}}$ is \textbf{bounded above/bounded below/bounded} if $A$ is such
			\end{enumerate} 
			
			\textit{Monotone Convergence Theorem For Sequences}: If $(a_n)_{n \in \mathbb{N}}$ is a real sequence which is monotone non-decreasing and bounded above; or monotone non-increasing and bounded below; then $(a_n)_{n \in \mathbb{N}}$ converges.
			
		\end{dfn}
		
		
		\begin{mdfn}[Cauchy Sequences]{mdfn:cauchyseq}{ }
			A real sequence $(a_n)_{n \in \mathbb{N}}$ is \textbf{Cauchy} if for every $\epsilon > 0$ there exists $N \in \mathbb{N}$ s.t. $\left|a_n - a_m\right| <\epsilon$ for all $n, m \ge N$.\\
			
		\end{mdfn}
		
		
		\begin{mthm}[The Cauchy Criterion]{thm:}{ }
			A real sequence is convergent if and only if it is Cauchy.
			
			\begin{proof}
				%TO-DO - PROOF
			\end{proof}
			
		\end{mthm}
		
		
		\begin{dfn}[Subsequences]{dfn:subseqs}{ }
			Let $(a_n)_{n \in \mathbb{N}}$ be a real sequence. A \textbf{subsequence} of $(a_n)_{n \in \mathbb{N}}$ is a real sequence of the form $(a_{n_k})_{k \in \mathbb{N}}$ where for each $k \in \mathbb{N}$ there exists $n_k \in \mathbb{N}$ such that $n_1 < n_2< ... < n_k < n_{k+1} < ...$. Ie, a subsequence is just some (possibly all) elements of a sequence taken in order.
			
		\end{dfn}
		
		
		\begin{thm}[Bolzano-Weierstrass Theorem]{thm:bolzweier}{ }
			Every bounded sequence of real numbers has a convergent subsequence.
			
			\begin{proof}
				%TO-DO - PROOF
			\end{proof}
			
		\end{thm}
		
		
		
		{\tiny \textbf{\underline{Series}}}
		
		
		
		\begin{mdfn}[Infinite Series]{dfn:infsers}{ }
			Let $(a_k)_{k \in N}$ be a real sequence. We form the \textbf{sequence of partial sums $(s_n)_{n \in \mathbb{N}}$} by $s_n := \sum_{k=1}^{n} a_k$ for all $n \in\mathbb{N}$.
			\begin{itemize}
				\item $\sum_{k=1}^{\infty}a_k$ \textbf{converges to $a$} if $\lim_{n\to \infty} s_n$ exists and $= a$, else \textbf{diverges}.
			\end{itemize} 
			
		\end{mdfn}
		
		
		\begin{dfn}[Absolute Convergence]{dfn:absconv}{ }
			Let $(a_n)_{n \in \mathbb{N}}$ be a real sequence.
			\begin{enumerate}[label=(\arabic*)]
				\item We say $\sum_{n=1}^{\infty} a_n$ \textbf{converges absolutely} if $\sum_{n=1}^{\infty} \left|a_n\right|$ converges
				\item We say $\sum_{n=1}^{\infty}a_n$ \textbf{converges conditionally} if it converges but $\sum_{n=1}^{\infty}\left|a_n\right|$ diverges.
			\end{enumerate}
			
			Absolute convergence implies convergence.
			\begin{proof}
				%TO-DO - PROOF
			\end{proof}
			
		\end{dfn}
		
		
		\begin{mthm}[The Cauchy Criterion For Series]{thm:cauchforseries}{ }
			Let $(a_k)_{k=1}^\infty$ be a real sequence. Then $\sum_{k=1}^{\infty}a_k$ converges iff for all $\epsilon > 0$ there exists $N \in \mathbb{N}$ such that $\left|\sum_{k=m+1}^{n}a_k\right| <\epsilon$ for all $n > m \ge N$.
			
		\end{mthm}
		
		
		\begin{thm}[The Riemann Rearrangement Theorem]{thm:riearrange}{ }
			Let $(a_n)_{n \in \mathbb{N}}$ be a real sequence and suppose $\sum_{n=1}^{\infty}a_n$ is conditionally convergent. Then for \textit{any} $s \in \mathbb{R}$, there exists a bijection $\sigma: \mathbb{N} \to \mathbb{N}$ such that the rearrangement $\sum_{n=1}^{\infty}a_{\sigma(n)}$ converges to $s$. 
			\vspace{1pt}
			\hrule
			\vspace{2pt}
			Note the RRT does not apply to absolutely convergent series. Indeed, for $\sum_{n=1}^{\infty}a_n$ absolutely convergent, and any bijection $\sigma: \mathbb{N} \to \mathbb{N}$, $\sum_{n=1}^{\infty}a_{\sigma(n)}$ is convergent and $\sum_{n=1}^{\infty}a_{\sigma(n)} = \sum_{n=1}^{\infty}a_n$
			
		\end{thm}
		
		
		
		{\tiny \textbf{\underline{Subsets Of The Real Line}}}
		
		
		
		\begin{dfn}[Open Sets]{dfn:opensets}{ }
			We say a set $U \subseteq \mathbb{R}$ is \textbf{open} if for all $x \in U$ there exists $r > 0$ such that $(x-r, x+r) \subseteq U$.
			
			\begin{itemize}
				\item The union of any collection of open sets is open
				\item The intersection of any finite collection of open sets is open
				\item Any nonempty open $U \subseteq \mathbb{R}$ is a finite or countable union of nonempty disjoint open intervals.
				\item If $f: \mathbb{R} \to \mathbb{R}$ is continuous and $\lambda \in \mathbb{R}$, then $\{x \in \mathbb{R} : f(x) > \lambda\}$ is open.
			\end{itemize}
			
		\end{dfn}
		
		
		\begin{dfn}[Closed Sets]{dfn:closedsets}{ }
			A set $F \subseteq \mathbb{R}$ is \textbf{closed} if its complement $\mathbb{R} \backslash F$ is open.
			
			\begin{itemize}
				\item The union of any finite collection of closed sets is closed
				\item The intersection of any collection of closed sets is closed
				\item If $f: \mathbb{R} to \mathbb{R}$ is continuous and $ \lambda \in \mathbb{R}$, then $\{x \in \mathbb{R} : f(x) \le \lambda\}$ is closed.
			\end{itemize}
			
		\end{dfn}
		
		
		\begin{dfn}[Sequential Closure]{dfn:seqclos}{ }
			$F \subseteq \mathbb{R}$ is \textbf{sequentially closed} if for every convergent sequence $(x_n)_{n \in \mathbb{N}}$ s.t. $x_n \in F$ for all $n \in \mathbb{N}$, the limit $x := \lim_{n\to \infty}x_n$ satisfies $x \in F$. \\
			$F \subseteq \mathbb{R}$ is closed if and only if it is sequentially closed.
			\begin{proof}
				%TO-DO - PROOF
			\end{proof}
			
		\end{dfn}
		
		
		\begin{dfn}[Level Sets]{dfn:lvlsets}{ }
			%TO-DO - ALL
			
		\end{dfn}
		
		
		\begin{dfn}[The Middle Thirds Cantor Set]{dfn:midthrdscantor}{ }
			The \textbf{Middle Thirds Cantor Set $C$} is a closed set defined by the following recursive process: 
			\begin{enumerate}[label=(\arabic*)]
				\item Divide $[0, 1]$ into thirds, remove the open middle third interval $(\frac{1}{3}, \frac{2}{3})$.
				\item Divide the remaining two intervals into thirds and remove the open middle thirds intervals $(\frac{1}{9}, \frac{2}{9})$ and $(\frac{7}{9}, \frac{8}{9})$.
				\item Repeat this process ad infinitum.
			\end{enumerate}
			At the $n^{th}$ stage of the construction we have $2^n$ closed intervals $(I_{n, k})_{k=1}^{2n}$ with $\lambda(I_{n, k}) = 3^{-n}$ for all $1 \le k \le 2^n$. \\
			Then $C := \bigcap_{n =1}^{\infty}C_n$ where $C_n :+ \bigcup_{k =1}^{2^n}I_{n, k}$ for $n \in \mathbb{N}$. Each $I_{n, k}$ is a closed interval so their finite union $C_n$ is closed, and $C$, the intersection of these closed $C_n$, is also closed.
			
		\end{dfn}
		
		
		
	%--------------------------------------------------------------------------------------%
	%%%%%%%%%%%%%%%%%%%%%%%%%%%%%%%%%%%%% SECTION 2 %%%%%%%%%%%%%%%%%%%%%%%%%%%%%%%%%%%%%%%%
	%--------------------------------------------------------------------------------------%
	
	
	{\footnotesize \textbf{\underline{Uniform Continuity, Uniform Convergence, Power Series}}}
	
	
	{\tiny\textbf{\underline{Uniform Continuity}}}
		
		
		
		\begin{mdfn}[Uniform Continuity]{dfn:unifcont}{ }
			Let $I \subseteq \mathbb{R}$ be an interval and $f: I \to \mathbb{R}$. $f$ is \textbf{uniformly continuous} if $\forall \epsilon > 0$, $\exists \delta > 0$ s.t. $\left|f(x) - f(y)\right| <\epsilon$ $\forall x, y \in I$ s.t. $\left|x-y\right| < \delta$. Note the key difference with regular continuity - here $\delta$ may only depend on $\epsilon$ only, not $x$.
			
		\end{mdfn}
		
		
		\begin{thm}[The Heine-Cantor Theorem]{thm:heincan}{ }
			If $a < b$ and $f: [a, b] \to \mathbb{R}$ is continuous, then $f$ is uniformly continuous.
			
		\end{thm}
		
		
		
	{\tiny \textbf{\underline{Uniform Convergence}}}	
		
		
		
		\begin{dfn}[Pointwise Convergence]{dfn:ptwiseconv}{ }
			Let $E \subseteq \mathbb{R}$ be nonempty. A sequence of functions $f_n: E \to \mathbb{R}$ \textbf{converges pointwise on $E$ to $f: E \to \mathbb{R}$} if for all $x\in E$ and $ \epsilon > 0$, there exists $N\in \mathbb{N}$ s.t. $\left|f_n(x) - f(x)\right| <\epsilon$ for all $n \ge N$.
			
		\end{dfn}
		
		
		\begin{thm}[Pointwise Convergence Theorems]{thm:ptwiseconv}{ }
			\begin{itemize}[label=*]
				\item The pointwise limit of bounded functions is not always bounded
				\item The ptwise limit of cont funcs is not always cont. $\exists f_n, f: [0,1] \to \mathbb{R}$ s.t.
				\begin{itemize}
					\item $f_n$ is continuous for all $n \in \mathbb{N}$
					\item $f_n \to f$ pointwise on $[0, 1]$ as $n \to \infty$
					\item $f$ is not continuous
					\item Example: $f_n(x) := 0$ if $ 0 \le x \le \frac{1}{2} - \frac{1}{2n}$; $nx - \frac{n-1}{2}$ if $ \frac{1}{2} - \frac{1}{2n} \le x \le \frac{1}{2}+ \frac{1}{2n}$; $1$ if $\frac{1}{2} + \frac{1}{2n} \le x \le 1$
				\end{itemize}
				\item The pointwise limit of differentiable functions is not always differentiable. $\exists f_n, f: (-1, 1) \to \mathbb{R}$ s.t.:
				\begin{itemize}
					\item $f_n$ is differentiable for all $N \in \mathbb{N}$
					\item $f_n \to f$ pointwise on $(-1, 1)$ as $n \to \infty$
					\item $f$ is not differentiable
					\item Example: Let $f_n(x) := \sqrt{x^2 + \frac{1}{n}}$, limit is not diffable at $0$.
				\end{itemize}
				\item The pointwise limit of derivatives is not always the derivative of the pointwise limit. $\exists f_n, f : (-1, 1) \to \mathbb{R}$ s.t.
				\begin{itemize}
					\item $f_n$ is diffable for all $n \in \mathbb{N}$ and $f$ is diffable
					\item $f_n \to f$ pointwise on $(-1, 1)$
					\item $f'(0) \ne \lim_{n\to \infty}f'_n(0)$, ie $\left(\lim_{n\to \infty f_n}\right)'(0) \ne \lim_{n\to \infty}f_n'(0)$ - the derivate and ptwise lim operations do not commute.
					\item Example: Let $f_n(x) = \frac{\sin(2 \pi n x)}{2 \pi n}, f'_n = \cos(2\pi n x), f = 0$
				\end{itemize}
			\end{itemize}
			
		\end{thm}
		
		
		\begin{dfn}[Graph and Vertical Neighbourhood]{dfn:graphnvertneigh}{ }
			Let $f: E \to \mathbb{R}$. The \textbf{graph of $f$} is the set $\upgamma(f) : = \left\{(f, f(x) : x\in E)\right\} \subset \mathbb{R}^2$.\\
			For $ \epsilon > 0$, the $e-$vertical neighbourhood of the graph of $f$ is the set $\mathscr{N}(f,\epsilon) := \left\{(x, y) : x\in E, \left|y-f(x)\right| <\epsilon\right\} \subset \mathbb{R}^2$.
			
		\end{dfn}
		
		
		\begin{mdfn}[Uniform Convergence]{mdfn:unifcon}{ }
			Let $E \subseteq \mathbb{R}$ be nonempty. A sequence of functions $f_n: E \to \mathbb{R}$ \textbf{converges uniformly on $E$ to $f:: E \to \mathbb{R}$} if for all $\epsilon > 0$ there exists $N\in \mathbb{N}$ s.t. $\left|f_n(x) - f(x)\right| <\epsilon$ for all $x\in E, n \ge N$.
			\begin{itemize}
				\item Uniform convergence implies pointwise convergence. \\
				For counterexample of converse use $f_n, f: (0, 1]  \to \mathbb{R}$. $f_n(x) := 1 $ if $0 < x < \frac{1}{n}$; $0$ if $\frac{1}{n} \le x \le 1$ and $f(x) :=0$.
			\end{itemize}
			
		\end{mdfn}
		
		
		\begin{thm}[Uniform Convergence Properties]{thm:unifconv}{ }
			Let $E \subseteq \mathbb{R}$ be nonempty, $f_n : E \to \mathbb{R}$ be a sequence of funcs, and $f: E \to \mathbb{R}$.
			\begin{itemize}
				\item The following are equivalent:
				\begin{enumerate}[label=(\arabic*)]
					\item $f_n \to f$ uniformly on $E$ as $n \to \infty$
					\item $\sup_{x\in E}\left|f_n(x) - f(x)\right| \to 0$ as $n \to \infty$
					\item For all $\epsilon > 0$ there exists $N\in \mathbb{N}$ s.t. $\upgamma(f_n) \subset \mathscr{N}(f,\epsilon)$ for all $n \ge N$.
				\end{enumerate}
				\item \textit{Sequence test for uniform convergence}: Let $f_n \to f$ uniformly on $E$ as $n \to \infty$. If $(x_n)_{n\in \mathbb{N}}$ is a real sequence with $x_n\in E$ for all $n\in \mathbb{N}$, then $\left|f_n(x_n) - f(x_n)\right| \to 0$ as $n \to \infty$.
			\end{itemize}
			
		\end{thm}
		
		
		
		\begin{dfn}[Uniform Cauchy]{dfn:unifcauch}{ }
			Let $E \subseteq \mathbb{R}$ be nonempty and $f_n : E \to \mathbb{R}$ be a sequence of functions. We say $(f_n)_{n\in \mathbb{N}}$ is \textbf{uniformly Cauchy on $E$} if for all $\epsilon > 0$ there exists $N\in \mathbb{N}$ s.t. $\left|f_n(x) - f_m(x)\right| < \epsilon$ for all $x\in E$ and $m, n \ge N$.
			
		\end{dfn}
		
		
		\begin{mthm}[The Uniform Cauchy Criterion]{thm:unifcauchcrit}{ }
			$(f_n)_{n\in \mathbb{N}}$ converges unifly on $E$ iff $(f_n)_{n\in \mathbb{N}}$ is unifly Cauchy on $E$.
			\begin{proof}
				%TO-DO? MAYBE? - PROOF
			\end{proof}
			
		\end{mthm}
		
		
		\begin{thm-s}[Uniform Limits Of Continuous/Diffable Functions]{thm:uniflimscontdiff}{ }
			\begin{itemize}
				\item Let $I \subseteq\mathbb{R}$ be an interval, $f_n, f : I \to \mathbb{R}$ s.t. $f_n \to f$ unifly on $I$ as $n \to \infty$. For all $x_0\in I$, if each $f_n$ is continuous at $x_0$, then $f$ is continuous at $x_0$.
				\begin{proof}
					%TO-DO - PROOF
				\end{proof}
				\item A unif lim of diffable funcs is not always diffable. $\exists f_n, f : (-1, 1) \to \mathbb{R}$ s.t.:
				\begin{itemize}
					\item $f_n$ is differentiable for all $n\in \mathbb{N}$
					\item $f_n \to f$ uniformly on $(-1, 1)$ as $n \to \infty$
					\item $f$ is not differentiable
					\item Example: consider $f_n(x) := \sqrt{x^2 + \frac{1}{n}}$
				\end{itemize}
				\item Let $I \subseteq \mathbb{R}$ be a bnded open interval, $f_n : I \to \mathbb{R}$ a seq of diffable funcs s.t.:
				\begin{enumerate}[label=(\arabic*)]
					\item There exists $x_0\in I$ s.t. $(f_n(x_0))_{n\in \mathbb{N}}$ converges
					\item There exists $g: I \to \mathbb{R}$ s.t. $f_n' \to g$ uniformly on $I$ as $n \to \infty$.
				\end{enumerate}
				Then there exists diffable $f: I \to \mathbb{R}$ s.t. $f_n \to f$ uniformly on $I$ as $n \to \infty$. Alternatively $(\lim_{n\to \infty}f_n)'(x) = \lim_{n\to \infty}f_n'(x)$ for all $x\in I$, ie differentiation commutes with taking the uniform limit under our conditions.
			\end{itemize}
			
		\end{thm-s}
		
		
		\begin{dfn}[Uniform Convergence Of Series Of Functions]{dfn:unficonvfuncseries}{ }
			Let $E \subseteq \mathbb{R}$ be nonempty and $f_n : E \to \mathbb{R}$ be a sequence of real funcs. Define the \textbf{sequence of partial sums} $s_n: E \to \mathbb{R}$ as $s_n(x) := \sum_{k=1}^{n}f_k(x)$ for all $x\in E, n\in \mathbb{N}$. Then:
			\begin{enumerate}[label=(\arabic*)]
				\item $\sum_{n=1}^{\infty}f_n$ \textbf{converges ptwise on $E$} if $(s_n)_{n\in \mathbb{N}}$ converges ptwise on $E$.
				\item $\sum_{n=1}^{\infty}f_n$ \textbf{converges unifly on $E$} if $(s_n)_{n\in \mathbb{N}}$ converges unifly on $E$
				\item $\sum_{n=1}^{\infty}f_n$ \textbf{converges absolutely (pointwise) on $E$} if $\sum_{n=1}^{\infty}\left|f_n(x)\right|$ converges $\forall x\in E$.
			\end{enumerate}
			
		\end{dfn}
		
		
		\begin{thm}[On Termwise Differentiation]{thm:trmwisediff}{ }
			Let $I \subset \mathbb{R}$ be a bounded open interval and suppose $f_n: I \to \mathbb{R}$ is a seq. of diffable funcs satisfying:
			\begin{enumerate}[label=(\arabic*)]
				\item There exists $x_0\in I$ such that $\sum_{n=1}^{\infty}f_n(x_0)$ converges
				\item $\sum_{n=1}^{\infty}f_n'$ converges uniformly on $I$.
			\end{enumerate}
			Then there exists diffable $f : I \to \mathbb{R}$ such that $\sum_{n=1}^{\infty}f_n$ converges uniformly to $f$ on $I$, and $f'(x) = \left(\sum_{n=1}^{\infty}f_n\right)'(x) = \sum_{n=1}^{\infty}f_n'(x)$ for all $x\in I$.
			
		\end{thm}
		
		
		\begin{thm}[Weierstrass M-Test]{thm:weierm}{ }
			Let $E \subseteq \mathbb{R}$ be $\ne \emptyset$, $f_n: E \to \mathbb{R}$, and $(M_n)_{n\in \mathbb{N}}$ a non-negative seq satisfying:
			\begin{enumerate}[label=(\arabic*)]
				\item $\left|f_n(x)\right| \le M_n$ for all $x\in E, n\in \mathbb{N}$
				\item $\sum_{n=1}^{\infty}M_n$ converges
			\end{enumerate}
			Then $\sum_{n=1}^{\infty}f_n$ converges absolutely and uniformly on $E$.
			\begin{proof}
				%TO-DO - PROOF
			\end{proof}
			
		\end{thm}
		
		
		\begin{thm}[Continuity Of Uniform Limits]{thm:uniflimscont}{ }
			Let $I \subseteq \mathbb{R}$ be an interval, $f_n, f : I \to \mathbb{R}$ s.t. $\sum_{n=1}^{\infty}f_n$ cnvges unifly to $f$ on $I$. For $x_0\in I$, if each $f_n$ is continuous at $x_0$, then $f$ is continuous at $x_0$.
			\begin{proof}
				%TO-DO - PROOF
			\end{proof}
			
		\end{thm}
		
		
		
		{\tiny \textbf{\underline{Power Series}}}
		
		
		
		\begin{dfn}[Power Series]{dfn:powseries}{ }
			Let $(a_n)_{n\in \mathbb{N}}$ be a real sequence and $c\in \mathbb{R}$. A \textbf{power series centred at $c$} is a series of the form $\sum_{n=0}^{\infty}a_n(x-c)^n$.
			
		\end{dfn}
		
		
		\begin{thm}[Radius Of Convergence]{thm:radofconv}{ }
			Let $(a_n)_{n\in \mathbb{N}}$ be a real seq and $c\in \mathbb{R}$. Then one of the following holds:
			\begin{enumerate}[label=(\arabic*)]
				\item $\sum_{n=0}^{\infty}a_n(x-c)^n$ converges only when $x = c$
				\item $\sum_{n=0}^{\infty}a_n(x-c)^n$ converges absolutely for all $x\in \mathbb{R}$
				\item There exists $R > 0, R\in \mathbb{R}$ such that $\sum_{n=0}^{\infty}a_n(x-c)^n$ converges absolutely when $\left|x - c\right| < R$ and diverges when $\left|x-c\right|> R$. Here $R = \sup\{r > 0 : (a_nr^n)_{n\in \mathbb{N}_0} \text{ is bounded }\}$
			\end{enumerate}
			The $R$ in case $(3)$ is the \textbf{radius of convergence} of $\sum_{n=0}^{\infty}a_n(x-c)^n$. In cases $(1)$ and $(2)$ we say $R=0$ and $R = \infty$ respectively.
			
		\end{thm}
		
		
		\begin{thm}[Continuity and Differentiability Of Power Series]{thm:contndiffpwrseries}{ }
			Let $\sum_{n=0}^{\infty}a_n(x-c)^n$ have radius of convergence $0 < R \le \infty$. Then $f(x) = \sum_{n=0}^{\infty}a_n(x-c)^n$ for $x\in (c-R, c+R)$ is differentiable on $(c-R, c+R)$; and $f'(x) = \sum_{n=1}^{\infty}na_n(x-c)^{n-1}$ for all $x\in (c-R, c+R)$.
			\begin{itemize}
				\item Given the above, $f(x) := \sum_{n=0}^{\infty}a_n(x-c)^n$  is \textbf{infinitely differentiable} on $(c-R, c+R)$, and $a_k = \frac{f^(k)(c)}{k!}$ for all $k\in \mathbb{N}_0$.
				\item For any $0 < r < R$, both of the above series converge absolutely and uniformly on $[c-r, c+r]$
				\item For any real sequence $(a_n)_{n\in \mathbb{N}_0}$ and $c\in \mathbb{R}$, the $\sum_{n=0}^{\infty}a_n(x-c)^n$ and $\sum_{n=1}^{\infty}na_n(x-c)^{n-1}$ have the same radius of convergence
			\end{itemize}
			
		\end{thm}
		
		
		\begin{mthm}[Uniqueness Of Power Series]{thm:uniqpwrsers}{ }
			Power series are unique - For $\sum_{n=0}^{\infty}a_n(x-c)^n$ and $\sum_{n=0}^{\infty}b_n(x-c)^n$ with RoCs $0 < R, S, \le \infty$, if $\exists 0 < r < \min{R, S}$ s.t. $\sum_{n=0}^{\infty} a_n(x-c)^n = \sum_{n=0}^{\infty}b_n (x-c)^n$ for all $x\in \mathbb{R}$ s.t. $\left|x-c\right| < r$, then $a_n = b_n$ for all $n\in \mathbb{N}$.
			\begin{proof}
				%TO-DO - PROOF
			\end{proof}
			
		\end{mthm}
		
		
		\begin{dfn}[The Exponential]{dfn:exp}{ }
			Consider $\sum_{n=0}^{\infty}\frac{x^n}{n!}$ with RoC $R = \infty$. $\exp(x) := e^x := \sum_{n=0}^{\infty}\frac{x^n}{n!}$.
			
		\end{dfn}
		
		
		\begin{dfn}[Taylor Series]{dfn:tayser}{ }
			Let $I \subseteq \mathbb{R}$ be a nonempty open interval, and $f: I \to \mathbb{R}$ be infinitely diffable. Then the \textbf{Taylor series of $f$ at $c\in I$} is the PS $\sum_{n=0}^{\infty}\frac{f^{(n)}(c)}{n!}(x-c)^n$.\\
			$f$ is \textbf{real analytic} at $x\in I$ if its TS converges in a neighbourhood of $x$, and to $f(x)$ at $x$.
			
		\end{dfn}
		
		
		
	%--------------------------------------------------------------------------------------%
	%%%%%%%%%%%%%%%%%%%%%%%%%%%%%%%%%%%%% SECTION 3 %%%%%%%%%%%%%%%%%%%%%%%%%%%%%%%%%%%%%%%%
	%--------------------------------------------------------------------------------------%
	
	
	{\footnotesize \textbf{\underline{Measure And Integration}}}
	
	
	{\tiny\textbf{\underline{The Lebesgue Measure}}}
		
		
		
		\begin{dfn}[The Lebesgue Outer Measure]{dfn:lebesoutmes}{ }
			For any $E \subseteq \mathbb{R}$, the \textbf{Lebesgue outer measure of $E$} is $m*(E) := \inf \left\{\sum_{j=1}^{\infty}\lambda(I_j) : E \subseteq\bigcup_{j=1}^\infty I_j, I_j \subseteq \mathbb{R} \text{ an open interval}\right\}$. Here the inf is taken over all countable coverings $(I_j)_{j\in \mathbb{N}}$ of $E$ by open intervals.
			\begin{itemize}
				\item[(1)] $m*(\emptyset) = 0$
				\item[(2)] If $E \subseteq F \subseteq \mathbb{R}$, $m*(E) \le m*(F)$
				\item[(3)] If $(E_j)_{j\in \mathbb{N}}$ is any seq of subsets of $\mathbb{R}$, $m*\left(\bigcup_{j=1}^{\infty} E_j\right) \le \sum_{j=1}^{\infty}m*(E_j)$
				\item If $\sum_{j=1}^{\infty}\lambda(I_j)$ diverges, we say it has value $\infty$. If there is no cover of $E$ s.t. $\sum_{j=1}^{\infty}\lambda(I_j) < \infty$, $m*(E) = \infty$.
				\item For any interval $J \subseteq\mathbb{R}$, $\lambda(J) = m*(J)$.
				\item Given $E \subseteq\mathbb{R}$, $\epsilon > 0$, $\exists$ open $U \subseteq \mathbb{R}$ s.t. $E \subseteq U$ and $m*(U) < m*(E) +\epsilon$.
			\end{itemize}
			
		\end{dfn}
		
		
		\begin{dfn}[Null Sets]{dfn:nullsets}{ }
			A set $E \subseteq \mathbb{R}$ is \textbf{null} or \textbf{has measure zero} if $m*(E) = 0$.
			\begin{itemize}
				\item Every countable subset of $R$ has measure zero.
				\begin{proof}
					%TO-DO - PROOF
				\end{proof}
			\end{itemize}
			
		\end{dfn}
		
		
		\begin{dfn}[Distance Between Sets]{dfn:dist}{ }
			Let $E, F \subseteq \mathbb{R}$ be nonempty and disjoint. Then dist$(E, F) := \inf\left\{\left|x-y\right|: x\in E, y\in F\right\}$.
			\begin{itemize}
				\item If $E, F \subseteq \mathbb{R}$ nonempty, disjoint, dist$(E, F) > 0$; $m*(E \cup F) = m*(E) + m*(F)$
				\item If $F< K \subseteq \mathbb{R}$ nonempty, closed, $F \cup K = \emptyset$, $K$ bounded; dist$(F, K) > 0$.
				\item If $\left(K_j\right)_{j\in \mathbb{N}} \subseteq\mathbb{R}$ a seq. of disjoint closed sets, $m*\left(\bigcup_{j=1}^\infty K_j\right) = \sum_{j=1}^{\infty}m*(K_j)$
			\end{itemize}
			
		\end{dfn}
		
		
		\begin{mdfn}[Lebesgue Measurable Sets]{mdfn:lebmesable}{ }
			$E \subseteq \mathbb{R}$ is \textbf{Lebesgue measurable} if for all $\epsilon > 0$, there exists open $U \subseteq \mathbb{R}$ s.t. $E \subseteq U$ and $m*(U \backslash E) <\epsilon$. The collection of all L-mesable subsets of $\mathbb{R}$ is denoted $\mathcal{M}$.
			
		\end{mdfn}
		
		
		\begin{thm}[Properties Of L-Measurable Sets]{thm:lmesprops}{ }
			\begin{itemize}
				\item Let $F \subseteq \mathbb{R}$ and $(E_j)_{j\in \mathbb{N}}$ be L-mesable. Then:
					\begin{enumerate}[label=(\arabic*)]
						\item The union $E := \bigcup_{j=1}^\infty E_j$ is L-mesable
						\item The complement $\mathbb{R} \backslash F$ is L-mesable
						\item The intersection $\bigcap_{j=1}^\infty E_j$ is L-mesable
					\end{enumerate}
				\item Let $E, F \subseteq \mathbb{R}$ be L-mesable. Then $F \backslash E$ is L-mesable
				\item For $E \subseteq \mathbb{R}$, $E$ is L-mesable iff $\forall \epsilon > 0$, $\exists$ closed $K \subseteq \mathbb{R}$ s.t. $K \subseteq E$ and $m*(E \backslash K) <\epsilon$.
			\end{itemize}
			
		\end{thm}
		
		
		
		
		
		
		
		
		
		
	\end{multicols*}
	
	
\end{document}
